\chapter{Project: Filtering Web Data for Language Modeling (CS336 Assignment 4)}\label{ch:proj4}

This project guide accompanies the data assignment of Stanford's CS336, \emph{Language Modeling from
Scratch}: take raw Common Crawl, turn it into text, filter, mask and deduplicate it, and show
that a model trained on your data reaches lower validation loss. The guide explains how the
assignment is set up and how to run it, maps each problem to the parts of this book that teach the
underlying ideas, and gives milestones, sanity checks and pitfalls. Like every guide in Part~VII,
it deliberately stops short of solutions (see the note at the start of \cref{ch:proj1}): the
point of the assignment is to make these decisions yourself, and the course asks that AI tools and
third-party code not do the problems for students.

\begin{figure}[t]
\centering
\resizebox{\linewidth}{!}{%
\begin{tikzpicture}[
  p/.style={draw=sysblue,fill=sysbluebg,rounded corners=2pt,minimum height=9mm,text width=23mm,align=center,font=\scriptsize},
  q/.style={p,fill=keygreenbg,draw=keygreen},
  r/.style={p,fill=napkinbg,draw=napkin},
  arr/.style={-{Stealth[length=2mm]},thick,black!55}]
\node[p] (look) {look\_at\_cc\\{\tiny read WARC/WET}};
\node[p,right=4mm of look] (ext) {extract\_text\\{\tiny bytes $\to$ text}};
\node[p,right=4mm of ext] (lang) {language\_\\identification};
\node[p,right=4mm of lang] (pii) {mask\_pii};
\node[p,right=4mm of pii] (harm) {harmful\_\\content};
\node[q,below=8mm of harm] (gop) {gopher\_quality\_\\filters};
\node[q,left=4mm of gop] (qc) {quality\_\\classifier};
\node[q,left=4mm of qc] (ex) {exact\_\\deduplication};
\node[q,left=4mm of ex] (mh) {minhash\_\\deduplication};
\node[r,below=8mm of mh] (filt) {filter\_data\\{\tiny parallel pipeline}};
\node[r,right=4mm of filt] (insp) {inspect\_\\filtered\_data};
\node[r,right=4mm of insp] (tok) {tokenize\_data\\{\tiny uint16 + EOT}};
\node[r,right=4mm of tok] (train) {train\_model\\{\tiny fixed recipe}};
\draw[arr] (look)--(ext); \draw[arr] (ext)--(lang); \draw[arr] (lang)--(pii); \draw[arr] (pii)--(harm);
\draw[arr] (harm)--(gop); \draw[arr] (gop)--(qc); \draw[arr] (qc)--(ex); \draw[arr] (ex)--(mh);
\draw[arr] (mh)--(filt); \draw[arr] (filt)--(insp); \draw[arr] (insp)--(tok); \draw[arr] (tok)--(train);
\draw[arr,dashed] (insp.north) to[out=90,in=-90] node[right,font=\tiny,text=srcgray] {iterate} (ex.south);
\node[font=\scriptsize,text=sysblue,anchor=south west] at (look.north west) {Part 1: primitives on one WARC/WET pair};
\node[font=\scriptsize,text=keygreen,anchor=south east] at (gop.north east) {Part 2: quality and dedup};
\node[font=\scriptsize,text=napkin,anchor=north west] at (filt.south west) {Part 3: build a dataset, train on it, measure};
\end{tikzpicture}}
\caption{The assignment's problems in order. Parts 1 and 2 build and test primitives; Part 3
composes them into a parallel pipeline over thousands of WET files, inspects the output, and is
judged by the validation loss of a model trained with a fixed recipe.}
\label{fig:proj4:map}
\end{figure}

% =====================================================================
\section{What the assignment asks and why}\label{sec:proj4:overview}

The assignment has three movements. First, you implement the per-document primitives
of \cref{ch:data}: HTML-to-text extraction with encoding detection, language identification,
PII masking, NSFW and toxicity classification, and Gopher-style quality rules, and you look hard at
what each one does to real Common Crawl pages. Second, you build a quality classifier from
positives you collect yourself (pages referenced by Wikipedia) and implement exact line deduplication
and MinHash-LSH fuzzy deduplication over files. Third, the open-ended part: filter a large collection
of English WET files into a training set, tokenise it, and train a fixed GPT-2-shaped model with the
course's training script. The score is that model's validation loss on a held-out benchmark (the
C4 100-domains subset of Paloma), so all of your leverage is in the \emph{data}; the architecture and
training procedure are frozen.

Most problems have two kinds of deliverable: code that passes provided tests, and a short written
analysis based on looking at samples. Take the analyses seriously. The habits they build (reading
20 random examples, tabulating classifier errors, picking thresholds from evidence) are exactly
what separates a working data pipeline from a plausible-looking one.

% =====================================================================
\section{Setting up}\label{sec:proj4:setup}

\subsection{Repository and environment}

The starter repository (\texttt{stanford-cs336/assignment4-data}) is organised as follows:

\begin{description}
  \item[\texttt{cs336\_basics/}] The staff implementation of the Assignment~1 model and a
  distributed (DDP) training script, slightly optimised with PyTorch-native pieces. You use it
  unchanged: the leaderboard requires the fixed recipe.
  \item[\texttt{cs336\_data/}] Your code. It starts nearly empty, apart from helpers for
  locating shared data and for building the English-filtered WET collection.
  \item[\texttt{tests/}] The test suite plus fixtures (an HTML page with its expected extraction,
  documents with line duplicates and with near-duplicate licences, high- and low-quality examples).
  Tests call functions in \texttt{tests/adapters.py}; your job is to make each adapter call
  \emph{your} implementation, keeping your real code in \texttt{cs336\_data}.
  \item[\texttt{scripts/}] Data download, training and generation entry points (Modal apps).
  \item[\texttt{configs/}] Hydra configs for training; you set the path to your tokenised data
  and your experiment-tracking entity and project.
\end{description}

Dependencies are managed with \texttt{uv}. The project pins the tools this assignment relies on:
FastWARC and warcio (record iteration), Resiliparse (encoding detection and extraction),
fastText (classifiers), mmh3 (MurmurHash), NLTK (tokenisation for rules), tldextract (domains),
submitit and Modal (parallel execution), and the training stack.

\begin{lstlisting}[style=shell]
git clone https://github.com/stanford-cs336/assignment4-data && cd assignment4-data
uv sync                              # creates .venv from uv.lock (Python >= 3.12)
uv run pytest -k test_extract        # run one family of tests while you work
uv run pytest                        # everything (expect failures until you finish)
\end{lstlisting}

\begin{pitfall}[Native wheels]
FastWARC, Resiliparse and fastText ship compiled extensions. If \code{uv sync} fails on your
platform (older glibc, Apple silicon with an unusual Python), check that \texttt{uv} is using the
managed Python the project requests, and read the specific build error before reaching for
workarounds. fastText's \code{predict} also rejects input containing newlines; strip them before
classifying.
\end{pitfall}

\subsection{Data}

The course provides a shared data volume containing an example WARC file with its WET counterpart,
the classifier binaries (fastText \code{lid.176.bin} for language ID and Dolma's Jigsaw-trained NSFW
and hate-speech classifiers), a gzipped list of external URLs extracted from English Wikipedia, a
large collection of English-filtered WET files (roughly 2{,}500 files, already filtered to an
English probability of at least 0.7 by the language-ID model), and the tokenised validation set. Students access it through Modal; the README
describes an offline download of just the files needed for the tests, and, for people outside the
course, how to point \texttt{cs336\_data/common.py} at a local directory, implement the English
check used when building the WET collection, and download fewer files.

\begin{napkinmath}[How much data, and what it costs to train]
The training recipe uses 8 GPUs, a per-device batch of 128 sequences of 512 tokens and 16{,}384 steps:
$8\times128\times512\times16384 = 2^{33}\approx8.6\times10^{9}$ tokens. A roughly 430M-parameter model
therefore costs about $6\times4.3\times10^{8}\times8.6\times10^{9}\approx2.2\times10^{19}$ FLOPs.
Finishing that in about two hours, as the handout reports, implies roughly
$3\times10^{15}$ FLOP/s across the node. On a single H100 at a realistic 40\% of its 989\,TFLOP/s
bf16 peak, the same run takes about 15 hours. Two consequences follow. First, if your filtered set
has fewer than 8.6B tokens, training will repeat it (\cref{sec:data:mixing}); know your epoch
count. Second, compare filters with \emph{short} runs before spending a full one.
\end{napkinmath}

\subsection{Compute and execution}

Filtering is CPU work and training is GPU work; plan them separately. The course runs both on
Modal: data scripts parallelise with \code{.map} over inputs, and training is launched with the
provided entry point and your tokenised file (\code{uv run modal run scripts/train.py --train-bin
<path>}). Off Modal, a process pool on a many-core machine (\code{concurrent.futures}) or a Slurm
array (\code{submitit}) does the filtering; the generic runner in \texttt{code/projects/a4} (a
resumable map with a manifest, counters and an ETA) and the pipeline of
\cref{sec:data:engineering} show the patterns. For training without eight large GPUs, see
\cref{sec:proj4:outside}.

\begin{lstlisting}[style=shell]
python code/projects/a4/env_check.py --data-root /path/to/shared-data
python code/projects/a4/inspect_samples.py --inputs example.warc.wet.gz \
    --n 25 --csv notes.csv
python code/projects/a4/shard_runner.py --fn my_filters:process_file \
    --inputs wet/*.warc.wet.gz --out-dir filtered/ --workers 32 --total-files 90000
\end{lstlisting}

% =====================================================================
\section{Problem-by-problem map}\label{sec:proj4:map}

\Cref{tab:proj4:map} lists every problem, what it asks in brief, and where this book develops the
ideas you need. The sections after the table add guidance specific to the assignment.

{\small
\begin{longtable}{@{}>{\raggedright\arraybackslash}p{0.21\linewidth}>{\raggedright\arraybackslash}p{0.41\linewidth}>{\raggedright\arraybackslash}p{0.32\linewidth}@{}}
\caption{CS336 Assignment 4 problems and where the background is taught.}\label{tab:proj4:map}\\
\toprule
\textbf{Problem} & \textbf{Asks you to\ldots} & \textbf{Background}\\
\midrule\endfirsthead
\toprule
\textbf{Problem} & \textbf{Asks you to\ldots} & \textbf{Background}\\
\midrule\endhead
look\_at\_cc & read raw WARC and WET records, judge the extracted text, annotate 25 documents & \cref{sec:data:warc}; \code{inspect\_samples.py}\\
extract\_text & decode bytes with encoding detection, extract text with Resiliparse, compare with WET & \cref{sec:data:extract}; Lab~11 tasks 1--2\\
language\_\allowbreak identification & wrap the fastText language-ID model, return a language and a score, choose a threshold from samples & \cref{sec:data:langid}; Lab~11 task 3\\
mask\_pii & mask emails, US phone numbers and IPv4 addresses with fixed placeholders, return counts, discuss errors & \cref{sec:data:pii}; Lab~11 task 4\\
harmful\_\allowbreak content & wrap the NSFW and toxic-speech classifiers, audit errors, choose thresholds & \cref{sec:data:pii}\\
gopher\_\allowbreak quality\_\allowbreak filters & implement (at least) a subset of Gopher rules, compare with your own judgement & \cref{sec:data:rules}, \cref{tab:data:gopher}; Lab~11 task 5\\
quality\_\allowbreak classifier & collect positives from Wikipedia-referenced URLs and negatives from CC, train a classifier, expose a label and a confidence & \cref{sec:data:quality}, \cref{fig:data:targetraw}; Lab~11 task 6\\
exact\_\allowbreak deduplication & two-pass line dedup across files with hashed counters & \cref{sec:data:exact}; Lab~11 task 7\\
minhash\_\allowbreak deduplication & normalise, shingle, MinHash, LSH bands, verify Jaccard, cluster, keep one per cluster, write files & \cref{sec:data:minhash,sec:data:lsh}; Lab~11 task 8\\
filter\_data & a parallel pipeline over the WET collection, with a per-filter breakdown and runtime estimates & \cref{sec:data:engineering}; Lab~11 task 11; \texttt{projects/a4}\\
inspect\_\allowbreak filtered\_\allowbreak data & read kept and removed samples, justify decisions, iterate & \cref{sec:data:engineering}\\
tokenize\_\allowbreak data & GPT-2 tokenise with an EOT after every document, write \code{uint16} & \cref{sec:text:pipeline}, \cref{sec:data:engineering}\\
train\_model & train the fixed model, report the best validation loss and learning curve & \cref{ch:pretrain}, \cref{sec:proj4:outside}\\
\bottomrule
\end{longtable}}

% =====================================================================
\section{Guidance for Part 1: per-document primitives}\label{sec:proj4:part1}

\paragraph{Look before you build.} Use \code{zcat ... | less} or the sample inspector, and keep
your 25 annotations in a CSV with columns for language, domain, page type and a keep/drop verdict.
That file becomes a tiny labelled evaluation set. Later you can measure every filter's precision and recall
on it instead of arguing from single examples.

\paragraph{Extraction.} The extractor wants a \code{str}, so the real work is decoding robustly:
try UTF-8, fall back to a detected encoding, never crash on a malformed page
(\cref{lst:data:extract} shows the cascade). When you compare your extraction with the WET text,
look at boilerplate (menus, footers), lists and tables, and whitespace. Neither extraction is
``correct''; describe what each one keeps that a model should or should not learn from.

\paragraph{Classifier wrappers.} For the fastText-based tasks (language, NSFW, toxicity), the
engineering is small but easy to get subtly wrong: load each model once per process, not once per call;
strip newlines from the input; map fastText's \texttt{\_\_label\_\_xx} strings to the identifiers
the tests expect; and return a score in $[0,1]$ that means what your label says. The written parts
ask for error analysis and thresholds: sample about 20 documents, label them yourself, tabulate
disagreements, and choose a threshold that makes the errors you can tolerate. A threshold justified
with numbers from your own sample is worth more than one copied from a paper.

\paragraph{PII.} Start from the formats listed in the problem, then write adversarial
examples (versions, prices, dates, ISBNs, international numbers) before running on real pages.
Report false positives and false negatives with examples. The downstream question, what a model
learns from placeholder-riddled text, has no single right answer; reason about it with
\cref{sec:data:pii}.

\paragraph{Gopher rules.} Implement the required subset exactly as specified (word-count bounds,
mean word length, ellipsis lines, alphabetic words). How you split text into words matters: NLTK's
tokenizer and \code{str.split} disagree on punctuation and change the mean word length. Record which
rule fired for every document; you will need those counts for \code{filter\_data}.

\begin{keyidea}[Every primitive gets a test you wrote yourself]
The provided tests are sanity checks, and some of them check only one or two examples. For each primitive, write three
or four tests of your own from the failure modes in \cref{ch:data}: an undeclared cp1252 page, a
short code-switched paragraph, a version number that looks like an IP address, a page that is all
bullets. These tests protect you when you start optimising for speed in Part~3.
\end{keyidea}

\subsection{How many samples is ``enough''?}\label{sec:proj4:samples}

Several problems ask you to estimate a filter's error rate from about 20 hand-labelled examples.
Treat such estimates as the rough signals they are. For a proportion $\hat p = k/n$, the Wilson 95\%
interval is
\[
\frac{\hat p + \frac{z^2}{2n} \pm z\sqrt{\frac{\hat p(1-\hat p)}{n} + \frac{z^2}{4n^2}}}{1+\frac{z^2}{n}},
\qquad z = 1.96.
\]
With $k=2$ errors in $n=20$ samples the interval is roughly $[3\%, 30\%]$; seeing \emph{zero}
errors in 20 still allows an error rate up to about 16\%. Pinning an error rate near 10\% to within
$\pm5$ points needs about $1.96^2\cdot0.1\cdot0.9/0.05^2\approx140$ labelled samples. Two practical
consequences follow. Report intervals, not just point estimates, in the writeup. And spend the
labelling budget where the decision is made: sample documents whose scores lie \emph{near} your
threshold, which is where a mistake would change what gets kept (\cref{ch:eval} covers the statistics in
more depth).

% =====================================================================
\section{Guidance for Part 2: the quality classifier and deduplication}\label{sec:proj4:part2}

\paragraph{Positives are a data-collection project.} The URL list is huge; subsample it. Fetching
pages yourself means dealing with timeouts, dead links, redirects and non-HTML responses. Use a
WARC-writing fetcher (the handout shows \code{wget} with \code{--warc-file} and a timeout), fetch in
parallel but politely, and expect a large fraction of the URLs to fail. Pass the fetched pages
through your own extraction, language, safety and rule filters before treating them as positives.
Otherwise the classifier learns the boilerplate of reputable sites.

\paragraph{Negatives and evaluation.} Draw negatives from the same stage of the pipeline at which
the classifier will be applied (for example, after language ID), balance the classes or use class weights,
and hold out a test split. Inspect the top- and bottom-scoring held-out documents: if the top
looks like navigation pages from news sites, your positives are contaminated. fastText's
supervised format is one example per line, \texttt{\_\_label\_\_<name> <text>}; word bigrams
(\code{wordNgrams=2}) and a few epochs are reasonable starting points. Keep the trained model under
version control, or at least its training-set hash.

\paragraph{Exact line dedup.} The interface is file-based: given a list of input files and an
output directory, write same-named files with the corpus-repeated lines removed. Two passes over
the files, a hash-keyed counter, and careful preservation of line endings are the whole design
(\cref{sec:data:exact}). Check your output on a toy corpus where you can enumerate the answer by
hand, and keep in mind the pitfall of \cref{sec:data:exact}.

\paragraph{MinHash dedup.} Normalise (lower-case, NFD, strip accents, remove punctuation,
collapse whitespace) before shingling, as the handout requires, and compute shingles in
\emph{words} with the requested $n$. Use a seeded hash family so runs are reproducible. Verify LSH
candidates with the true Jaccard similarity, cluster transitively, and keep exactly one document
per cluster. Test with the provided licence fixtures, then with your own: two documents differing
in one word should merge, and two unrelated documents that share a long boilerplate footer should not
(if they do, your threshold or shingle size is wrong). \Cref{fig:data:scurve} and
\cref{eq:data:scurve} let you reason about which $(b,r)$ split makes a given threshold work before
you run anything.

\begin{pitfall}[Determinism and memory]
Python's built-in \code{hash()} is randomised per process: never use it for anything that must
agree across workers or runs. Use MurmurHash (\code{mmh3}) or a fixed-digest hash. For large inputs,
keep counters keyed by 64-bit integers rather than strings, and stream files line by line rather
than reading them into memory.
\end{pitfall}

% =====================================================================
\section{Guidance for Part 3: building the dataset}\label{sec:proj4:part3}

\paragraph{Design the pipeline on paper first.} Draw your funnel (\cref{fig:data:funnel}): which
stages run per file, which are global, in what order, and with which thresholds. Order stages by
cost and selectivity (\cref{sec:data:engineering}). Decide what you will log for every stage
(documents and characters in and out, time); the problem requires a breakdown of what each
filter removed, and you will want it anyway to debug.

\paragraph{Run it as a resumable parallel job.} Process each WET file independently in a process
pool or on Modal, write outputs atomically, and keep a manifest so that a crash costs one file, not
the run. Load classifiers once per worker. Time a few files first and extrapolate to the whole
collection and to a full crawl (the problem asks for both). \Cref{sec:data:engineering} has the
napkin method. Global deduplication then runs over the filtered outputs, which are much smaller.

\paragraph{The validation target shapes the data.} Your goal is low perplexity on a specific held-out
distribution: documents from the most common web domains in C4. The handout permits using the
validation data to \emph{design} filters or classifiers, for example as a target set $T$ in the sense of
\cref{fig:data:targetraw}, but forbids copying any validation text into training. Treat that rule
as a hard constraint and enforce it mechanically: run a long-$n$-gram decontamination pass
(\cref{sec:data:decontam}) of your final training set against the validation text, and report the
result. Think also about quantity versus quality (\cref{sec:data:quality}): the fixed recipe consumes
about 8.6B tokens, so aggressive filtering trades quality for repetition.

\paragraph{Cluster execution without Modal.} On a Slurm cluster, \code{submitit} turns a Python
function into a job array with one task per chunk of files. The pattern below is generic: the
function you submit is your own per-file processing function, run through a resumable wrapper like
\texttt{projects/a4/shard\_runner.py}.

\begin{lstlisting}
import submitit
files = sorted(glob.glob("/data/english-wet/*.warc.wet.gz"))
chunks = [files[i::256] for i in range(256)]           # 256 array tasks
ex = submitit.AutoExecutor(folder="logs/filter_%j")
ex.update_parameters(timeout_min=240, cpus_per_task=16, mem_gb=64,
                     slurm_array_parallelism=64)
jobs = ex.map_array(process_chunk, chunks, [OUT_DIR] * len(chunks))
counters = [j.result() for j in jobs]   # blocks; a re-run skips finished files
\end{lstlisting}

\noindent Inside each task, use a process pool over that task's files with
\code{len(os.sched\_getaffinity(0))} workers. Keep the per-file function idempotent, so that
re-submitting the whole array after a failure redoes only the missing files.

\begin{figure}[t]
\centering
\begin{tikzpicture}
\begin{axis}[width=0.72\linewidth,height=4.8cm,xmode=log,xlabel={cost per evaluation (seconds, log scale)},
  ylabel={fidelity to final metric},xmin=1,xmax=2e5,ymin=0,ymax=1.1,yticklabels={},
  label style={font=\small},tick label style={font=\scriptsize}]
\addplot[only marks,mark=*,mark size=2.5pt,sysblue] coordinates {(5,0.2) (300,0.35) (4000,0.75) (60000,1.0)};
\addplot[sysblue!40,thick] coordinates {(5,0.2) (300,0.35) (4000,0.75) (60000,1.0)};
\node[font=\scriptsize,anchor=south west] at (axis cs:5,0.22) {funnel + samples};
\node[font=\scriptsize,anchor=north west] at (axis cs:300,0.33) {decontamination, token count};
\node[font=\scriptsize,anchor=north west] at (axis cs:4600,0.73) {short proxy training run};
\node[font=\scriptsize,anchor=south east] at (axis cs:55000,1.01) {full fixed-recipe run};
\end{axis}
\end{tikzpicture}
\caption{The evaluation ladder for data decisions (schematic). Run the cheap rungs on every
pipeline variant and the expensive rungs only on the survivors. Costs are orders of magnitude, not
measurements.}
\label{fig:proj4:ladder}
\end{figure}

\paragraph{Inspect and iterate.} Read random kept and removed documents (the problem asks for five
of each with justifications). Then change one thing at a time and record what changed: a threshold,
a stage order, a new rule. An experiment log as simple as a table of (pipeline version, tokens kept,
per-filter counts, short-run validation loss) is what makes the final writeup straightforward.

\paragraph{Tokenise and train.} Encode with the GPT-2 tokenizer, append the end-of-text id after every
document, and write a single flat \code{uint16} array, the format the training script memory-maps
(\cref{sec:text:pipeline}). Parallelise tokenisation and avoid materialising one giant Python list.
Check the file before launching training: its size in bytes divided by two is the token count, and decoding
a random window should give readable text with EOT separators. Then train with the provided script,
log to your experiment tracker, and report the best validation loss with the learning curve.

\begin{systemslens}[Make every data decision cheap to evaluate]
A full training run is the most expensive evaluation you have. Build a ladder of cheaper proxies.
First, the per-filter funnel and samples (seconds). Second, the decontamination check (minutes). Third,
a short training run, for example 5--10\% of the steps or a smaller model, compared across pipeline variants with the
same seed and token budget (an hour). Last, the full run, only for the one or two best candidates.
\Cref{fig:proj4:ladder} sketches the trade-off. This is the data-ablation methodology of DCLM and
FineWeb (\cref{sec:data:engineering}) at the scale of one student; the scale caveat of \cref{sec:data:quality} still applies.
\end{systemslens}

% =====================================================================
\section{Background: selecting data toward a target distribution}\label{sec:proj4:selection}

Part~3 is a \emph{targeted} data-selection problem: choose, from a large raw pool $R$, the training
set that minimises loss on a specific evaluation distribution $T$. \Cref{sec:data:quality}
described quality classifiers as one instance of ``find documents like $T$''. This section adds the
general framework, so that you can reason about the design space rather than try filters at random.

\paragraph{Selection as importance resampling.} If training on $T$ itself were possible we would
do it, so the ideal is to sample from $R$ with weights proportional to the density ratio
\begin{equation}
w(x) = \frac{p_T(x)}{p_R(x)}.
\label{eq:proj4:dsir}
\end{equation}
Resampling $R$ with probability proportional to $w(x)$ produces a sample distributed like $T$
restricted to the support of $R$. Direct density estimation of text is hopeless, but cheap
\emph{feature-space} estimates work well: DSIR (Xie et al., 2023) fits hashed $n$-gram bag-of-words models
$\hat p_T$ and $\hat p_R$ (the same hashing trick as \cref{eq:data:hashed}) and resamples with
$\log w(x) = \sum_{f\in F(x)}\big(\log\hat p_T(f) - \log\hat p_R(f)\big)$, with Gumbel top-$k$ sampling
for efficiency. Two other estimators of the same ratio are already familiar. A calibrated classifier
gives $\frac{p(T\mid x)}{p(R\mid x)} \propto w(x)$, so thresholding its score is a hard
version of importance resampling. Perplexity under an LM trained on $T$, compared with perplexity under an
LM trained on $R$, gives the generative version (Moore and Lewis, 2010).

\begin{figure}[t]
\centering
\begin{tikzpicture}
\begin{axis}[width=0.78\linewidth,height=4.6cm,xlabel={feature space (schematic)},ylabel={density},
  xmin=-4,xmax=6,ymin=0,ymax=0.62,yticklabels={},xticklabels={},
  legend style={font=\scriptsize,at={(0.98,0.98)},anchor=north east},
  label style={font=\small},axis lines=left]
\addplot[domain=-4:6,samples=150,thick,black!50,fill=black!6] {0.5*exp(-(x-0)^2/2)/sqrt(2*pi) + 0.5*exp(-(x-2.5)^2/4)/sqrt(4*pi)} \closedcycle;
\addplot[domain=-4:6,samples=150,very thick,keygreen] {exp(-(x-2)^2/(2*0.36))/sqrt(2*pi*0.36)*0.55};
\addplot[domain=-4:6,samples=150,thick,dashed,sysblue] {min(0.6, 0.12*(exp(-(x-2)^2/(2*0.36))/sqrt(2*pi*0.36)*0.55)/(0.5*exp(-(x-0)^2/2)/sqrt(2*pi) + 0.5*exp(-(x-2.5)^2/4)/sqrt(4*pi)))};
\legend{raw pool $p_R$,target $p_T$,weight $w\propto p_T/p_R$ (scaled)}
\end{axis}
\end{tikzpicture}
\caption{Targeted selection as importance resampling. Documents are kept in proportion to how
much more likely they are under the target than under the raw pool. Weights explode where $p_R$
is small, so practical methods clip or threshold them; a classifier score is a monotone transform of
the same ratio.}
\label{fig:proj4:dsir}
\end{figure}

\paragraph{The questions that decide the design.} Before choosing a method, answer four questions
with data. \textbf{What is in $T$?} Read validation samples; tabulate domains, genres and lengths.
\textbf{Where is $T$ in $R$?} Score a sample of $R$ with your candidate estimator and read the
top-scoring documents. \textbf{How many tokens survive at each threshold?} Draw the curve of
tokens kept against threshold and mark the 8.6B-token budget on it. \textbf{Is the estimator
leaking?} Any method that looks at $T$ can memorise it; decontamination (\cref{sec:data:decontam})
is the mechanical guard. Diversity is the usual casualty of aggressive selection, so watch the
domain and length distributions of what you keep, not just the score.

\begin{pitfall}[Selection bias toward the estimator's features]
A hashed-unigram estimator rewards vocabulary overlap. It will happily select keyword-stuffed pages
that contain the target's frequent words. Combine targeted selection with the generic quality and
dedup stages, and check the top-scoring documents by eye before committing.
\end{pitfall}

% =====================================================================
\section{How the training side consumes your data}\label{sec:proj4:training}

Reading the provided training code before producing data saves wasted runs. The script memory-maps
your \code{uint16} file and draws random windows of the context length (the same random-window loader
as in \cref{ch:optim} and \cref{sec:text:pipeline}), so document order in the file does not matter,
but every token is a potential window start: a long run of junk anywhere in the file will be sampled.
Validation uses the provided pre-tokenised file with the same loader, periodically, and logs to your
experiment tracker. The model, optimiser, schedule and batch size come from the config and are
part of the fixed recipe. Your only input is the training file.

Three checks catch most integration bugs before a long run:

\begin{lstlisting}
import numpy as np
from transformers import AutoTokenizer
ids = np.memmap("my_data.bin", dtype=np.uint16, mode="r")
tok = AutoTokenizer.from_pretrained("gpt2")
print(f"{len(ids):,} tokens  ({len(ids) * 2 / 1e9:.2f} GB)")
assert ids.max() < tok.vocab_size                       # 50257 for GPT-2
eot = (ids == tok.eos_token_id).sum()
print(f"{eot:,} documents, mean {len(ids) / max(eot, 1):.0f} tokens/doc")
s = np.random.randint(0, len(ids) - 1024)
print(tok.decode(ids[s:s + 1024]))                       # read it!
\end{lstlisting}

\noindent The first check catches \code{int64} arrays written by mistake (four times the size and
garbage when read as \code{uint16}). The second catches missing end-of-text separators. The third
catches mojibake, leftover markup and placeholder floods. Comparing the mean document length with
your pipeline's counters closes the loop.

% =====================================================================
\section{Tracking data experiments}\label{sec:proj4:tracking}

Data work produces many artefacts: filtered shards, classifiers, thresholds, token files and
training runs. Months later you must be able to say which pipeline produced which loss. A light
discipline is enough:

\begin{itemize}
  \item \textbf{Config-driven pipeline.} Keep every threshold, stage order, model path and seed in one
        versioned config (YAML or a dataclass), and write it next to the output directory.
  \item \textbf{Content-addressed outputs.} Name dataset directories by a short hash of
        (config, code version, input file list), so that two different pipelines can never overwrite each other.
  \item \textbf{Manifests and counters.} Store per-file counters (the runner in \texttt{projects/a4}
        writes \code{\_manifest.jsonl}) and aggregate them into the funnel for the writeup.
  \item \textbf{Linked training runs.} Log the dataset hash, its token count and its funnel as run
        metadata in your tracker, so that every learning curve points back to its data.
  \item \textbf{Saved samples.} Store a fixed random sample of kept and removed documents for every
        dataset version; reviewing those samples is the qualitative half of every comparison.
\end{itemize}

\begin{table}[t]
\centering\small
\begin{tabularx}{\linewidth}{@{}>{\raggedright\arraybackslash}p{0.27\linewidth}>{\raggedright\arraybackslash}X>{\raggedright\arraybackslash}p{0.27\linewidth}@{}}
\toprule
\textbf{Symptom} & \textbf{Likely cause} & \textbf{Check}\\
\midrule
Pipeline keeps almost nothing & a threshold on the wrong scale (score of the wrong label), or a rule applied to raw HTML instead of text & funnel per stage on one file; print scores of 20 good pages\\
Kept text full of menus and footers & extraction keeps boilerplate; line dedup not applied or applied per file instead of across files & read 10 kept docs; count repeated lines across files\\
Workers slow down over time, then die & a model loaded per call; memory growth from accumulating results & time per file and RSS per worker; load once per process\\
Results differ between runs & built-in \code{hash()}, unseeded sampling, nondeterministic survivor choice & rerun on one file and diff outputs\\
Token file twice or four times too large & wrong dtype (\code{int32}/\code{int64}) & bytes / tokens should equal 2\\
Training loss fine, validation poor & distribution mismatch with the target, or too few unique tokens (heavy repetition) & tokens kept vs.\ 8.6B; domain mix of kept data\\
Validation loss suspiciously low & validation text leaked into training & long-$n$-gram decontamination report\\
\bottomrule
\end{tabularx}
\caption{Common failure modes in the data assignment and how to diagnose them.}
\label{tab:proj4:failures}
\end{table}

% =====================================================================
\section{Milestones, sanity checks and the writeup}\label{sec:proj4:milestones}

\begin{enumerate}
  \item \textbf{Environment}: \code{uv run pytest} collects all tests; the shared data or your local
        copy is reachable; \code{env\_check.py} is green.
  \item \textbf{Primitives}: each family of tests passes (\code{uv run pytest -k <name>}), plus your own
        adversarial tests; every written answer is backed by a small table from real samples.
  \item \textbf{Quality classifier}: a held-out accuracy or AUC you can quote, and top and bottom examples
        that look right.
  \item \textbf{Dedup}: the toy-corpus tests pass; on one real WET file, report how many lines and
        documents each dedup removes, and read some removed examples.
  \item \textbf{Pipeline}: one WET file end to end, with the funnel printed; then ten files in parallel with
        speedup measured; then everything, resumable.
  \item \textbf{Dataset}: the token count, the decontamination report, and the per-filter breakdown.
  \item \textbf{Model}: the learning curve and best validation loss, and if possible a short-run
        comparison of two pipeline variants.
\end{enumerate}

\paragraph{Sanity checks.} The number of documents in equals the number kept plus the number removed for every stage (counters
that do not add up mean a bug); \cref{tab:proj4:failures} lists symptoms and checks. Rerunning the pipeline produces byte-identical outputs (determinism).
The validation loss of a model trained on unfiltered English WET text is a useful baseline: your
filters should beat it, and if they do not, look at what you removed.

\paragraph{The writeup.} Organise it by problem, and for each problem state what you did, the evidence
(tables of samples, counts, plots), and the decision you took because of it. For Part~3 include the
funnel figure, the runtime and its extrapolation, the inspected examples, the iterations you tried,
and the final learning curve. Being honest about failed ideas is part of good data work.

% =====================================================================
\section{Doing this project outside Stanford}\label{sec:proj4:outside}

The primitives and the pipeline need only CPUs: a laptop handles Parts 1--2, and a single many-core
workstation can filter a few hundred WET files in hours. You can download Common Crawl segments yourself (the
handout shows the URL pattern) and the public classifier binaries; build a smaller
English-filtered collection with the repository's helper and fewer files.

Training is the expensive step. Three scaled-down options, in increasing fidelity:

\begin{itemize}
  \item Train the Lab~10 model (\cref{ch:pretrain}) or a GPT-2-small-shaped model on one GPU for a
        fixed, small token budget, and compare pipeline variants \emph{relative} to each other on the same
        validation file. This is the right tool for choosing between data variants.
  \item Run the course's training script on whatever GPUs you have with fewer steps or a smaller model. The
        script supports single-GPU and multi-GPU DDP launches (\cref{ch:distributed}).
  \item Rent a multi-GPU node for the final run only, after the cheaper rungs of the ladder have
        selected the data.
\end{itemize}

Whichever you choose, write down the budget (tokens, steps, model size) next to every reported
loss, because absolute numbers are comparable only within a fixed recipe.

% =====================================================================
\section{Beyond the assignment}\label{sec:proj4:beyond}

Once the required pipeline works, several directions are worth exploring as stretch work. They
connect to research and production practice:

\begin{description}
  \item[Semantic deduplication.] MinHash catches lexical near-duplicates; embedding-based methods
  (for example, clustering sentence embeddings and removing points that are close within a cluster,
  as in SemDeDup) also catch paraphrases. The cost is an embedding model forward pass per document,
  which changes the systems picture from CPU-bound to GPU-bound (\cref{sec:data:quality}'s cost
  napkin applies).
  \item[Small-LM perplexity filters.] Replace $n$-gram estimators with a small transformer trained on
  the target set. Perplexity-based pruning can work well but is sensitive to which reference model you use.
  \item[Rephrasing instead of discarding.] Use an instruction-tuned model to rewrite borderline pages
  into clean prose (\cref{sec:data:synthetic}). Measure the effect at a fixed token budget, and audit
  for factual drift.
  \item[Mixture search.] Split your kept data into buckets (by domain, classifier score or length)
  and search the mixture weights with small runs, respecting epochs (\cref{sec:data:mixing}).
  \item[Throughput engineering.] Profile the pipeline (\cref{ch:perf} for the general methodology), move
  the dominant stage to a compiled implementation, and measure parallel efficiency against core count
  (\cref{sec:data:engineering}).
\end{description}

% =====================================================================
\section{Checklist}\label{sec:proj4:checklist}

\begin{itemize}
  \item[$\square$] Environment builds; tests collect; data reachable; \code{env\_check.py} green.
  \item[$\square$] 25 annotated samples in a CSV; every written answer backed by a table or counts.
  \item[$\square$] All primitives pass the provided tests \emph{and} your own adversarial tests.
  \item[$\square$] Quality classifier: positives filtered, held-out metric, top/bottom examples read.
  \item[$\square$] Dedup: toy-corpus checks pass; deterministic hashing; one survivor per cluster.
  \item[$\square$] Pipeline: idempotent per-file jobs, manifest, per-stage counters that add up.
  \item[$\square$] Runtime measured on a sample and extrapolated to the full collection and a full crawl.
  \item[$\square$] Token file validated (dtype, EOT count, readable windows); decontamination reported.
  \item[$\square$] Proxy runs compare pipeline variants at a fixed budget before the full run.
  \item[$\square$] Writeup: funnel, samples, iterations, learning curve, budget stated next to every number.
\end{itemize}

\begin{deepdive}[Where to go further]
The open data pipelines are the best next reading: DataTrove (the FineWeb pipeline), Dolma's
toolkit, and DCLM's processing code all implement the stages of \cref{ch:data} at scale with
sharding, MinHash dedup and classifier filtering. Reading how they shard work and handle failures
is as instructive as reading how they filter. See \cref{app:refs}.
\end{deepdive}
